\documentclass[10pt,a4paper]{amsart}
\usepackage{amsmath,amssymb,amsthm,graphicx}
\usepackage{verbatim}
\newenvironment{reference}{\comment}{\endcomment}
%\newenvironment{reference}{}{}
\newenvironment{slogan}{\comment}{\endcomment}
\newenvironment{history}{\comment}{\endcomment}

% For commutative diagrams we use Xy-pic
\usepackage[all]{xy}

% We use 2cell for 2-commutative diagrams.
\xyoption{2cell}
% Theorem environments.
%
\theoremstyle{plain}
\newtheorem{theorem}[subsection]{Theorem}
\newtheorem{proposition}[subsection]{Proposition}
\newtheorem{lemma}[subsection]{Lemma}

\theoremstyle{definition}
\newtheorem{definition}[subsection]{Definition}
\newtheorem{example}[subsection]{Example}
\newtheorem{exercise}[subsection]{Exercise}
\newtheorem{situation}[subsection]{Situation}

\theoremstyle{remark}
\newtheorem{remark}[subsection]{Remark}
\newtheorem{remarks}[subsection]{Remarks}

\numberwithin{equation}{subsection}

% Macros
%
\def\lim{\mathop{\mathrm{lim}}\nolimits}
\def\colim{\mathop{\mathrm{colim}}\nolimits}
\def\Spec{\mathop{\mathrm{Spec}}}
\def\Hom{\mathop{\mathrm{Hom}}\nolimits}
\def\Ext{\mathop{\mathrm{Ext}}\nolimits}
\def\SheafHom{\mathop{\mathcal{H}\!\mathit{om}}\nolimits}
\def\SheafExt{\mathop{\mathcal{E}\!\mathit{xt}}\nolimits}
\def\Sch{\mathit{Sch}}
\def\Mor{\mathop{\mathrm{Mor}}\nolimits}
\def\Ob{\mathop{\mathrm{Ob}}\nolimits}
\def\Sh{\mathop{\mathit{Sh}}\nolimits}
\def\NL{\mathop{N\!L}\nolimits}
\def\CH{\mathop{\mathrm{CH}}\nolimits}
\def\proetale{{pro\text{-}\acute{e}tale}}
\def\etale{{\acute{e}tale}}
\def\QCoh{\mathit{QCoh}}
\def\Ker{\mathop{\mathrm{Ker}}}
\def\Im{\mathop{\mathrm{Im}}}
\def\Coker{\mathop{\mathrm{Coker}}}
\def\Coim{\mathop{\mathrm{Coim}}}

% Boxtimes
%
\DeclareMathSymbol{\boxtimes}{\mathbin}{AMSa}{"02}

%
% Macros for moduli stacks/spaces
%
\def\QCohstack{\mathcal{QC}\!\mathit{oh}}
\def\Cohstack{\mathcal{C}\!\mathit{oh}}
\def\Spacesstack{\mathcal{S}\!\mathit{paces}}
\def\Quotfunctor{\mathrm{Quot}}
\def\Hilbfunctor{\mathrm{Hilb}}
\def\Curvesstack{\mathcal{C}\!\mathit{urves}}
\def\Polarizedstack{\mathcal{P}\!\mathit{olarized}}
\def\Complexesstack{\mathcal{C}\!\mathit{omplexes}}
% \Pic is the operator that assigns to X its picard group, usage \Pic(X)
% \Picardstack_{X/B} denotes the Picard stack of X over B
% \Picardfunctor_{X/B} denotes the Picard functor of X over B
\def\Pic{\mathop{\mathrm{Pic}}\nolimits}
\def\Picardstack{\mathcal{P}\!\mathit{ic}}
\def\Picardfunctor{\mathrm{Pic}}
\def\Deformationcategory{\mathcal{D}\!\mathit{ef}}
\usepackage{cleveref}
\crefname{originalproposition}{original proposition}{original propositions}
\Crefname{originalproposition}{Original proposition}{Original propositions}
\makeatletter
\newlabel{proposition-loc-qcoh-flat-base-change}{{0771}{source}{}{sourcepointer-1}{}}
\newlabel{proposition-loc-qcoh-flat-base-change@cref}{{[originalproposition][0][]0771}{[1][0][]source}{}{}{}}
\makeatother
\crefname{originallemma}{original lemma}{original lemmas}
\Crefname{originallemma}{Original lemma}{Original lemmas}
\makeatletter
\newlabel{lemma-adjoint}{{0778}{source}{}{sourcepointer-2}{}}
\newlabel{lemma-adjoint@cref}{{[originallemma][0][]0778}{[1][0][]source}{}{}{}}
\makeatother
\crefname{originalreference}{original reference}{original references}
\Crefname{originalreference}{Original reference}{Original references}
\makeatletter
\newlabel{homology-definition-cohomological-delta-functor}{{010Q}{source}{}{sourcepointer-3}{}}
\newlabel{homology-definition-cohomological-delta-functor@cref}{{[originalreference][0][]010Q}{[1][0][]source}{}{}{}}
\makeatother
\crefname{originallemma}{original lemma}{original lemmas}
\Crefname{originallemma}{Original lemma}{Original lemmas}
\makeatletter
\newlabel{lemma-exact-sequence-quasi-coherent-parasitic-cohomology}{{0776}{source}{}{sourcepointer-4}{}}
\newlabel{lemma-exact-sequence-quasi-coherent-parasitic-cohomology@cref}{{[originallemma][0][]0776}{[1][0][]source}{}{}{}}
\makeatother
\crefname{originallemma}{original lemma}{original lemmas}
\Crefname{originallemma}{Original lemma}{Original lemmas}
\makeatletter
\newlabel{lemma-pushforward-parasitic}{{0775}{source}{}{sourcepointer-5}{}}
\newlabel{lemma-pushforward-parasitic@cref}{{[originallemma][0][]0775}{[1][0][]source}{}{}{}}
\makeatother
\crefname{originallemma}{original lemma}{original lemmas}
\Crefname{originallemma}{Original lemma}{Original lemmas}
\makeatletter
\newlabel{lemma-adjoint-kernel-parasitic}{{0779}{source}{}{sourcepointer-6}{}}
\newlabel{lemma-adjoint-kernel-parasitic@cref}{{[originallemma][0][]0779}{[1][0][]source}{}{}{}}
\makeatother
\title{The \$f\_\{\textbackslash{}QCoh, *\}\$ adjoint and its cohomological \$\textbackslash{}delta\$-functor structure}
\author{The Stacks Project Authors}
\date{stacks-project:stacks-cohomology.tex:local-snapshot-20261012; GNU FDL 1.2 or later}
\begin{document}
\maketitle
\noindent\textit{Editorial note.} All delivered mathematical text is the author's own, copied verbatim from the pinned Stacks Project chapter file stacks-cohomology.tex; no statement or proof step was rewritten, reordered or removed.\par\medskip
\noindent\textit{Editorial note.} The counted span cites items from Homology plus five same-chapter results, recorded in assembly.external\_references with canonical tags as navigation pointers.\par\medskip
\noindent\textit{Allowed standard prerequisites (editorial).}\begin{itemize}
\item Right adjoint \$Q\$ to the inclusion of \$\textbackslash{}QCoh\$. Stacks Project, Cohomology of Algebraic Stacks, Lemma lemma-adjoint. Supplies the second and third adjunction steps and exactness of \$Q\$.
\item Summary of LQCoh\textasciicircum{}\{fbc\} properties. Stacks Project, Cohomology of Algebraic Stacks, Proposition proposition-loc-qcoh-flat-base-change. Supplies membership of direct images and topology-independence.
\item Exactness of quasi-coherent sequences up to parasitic modules. Stacks Project, Cohomology of Algebraic Stacks, Lemma lemma-exact-sequence-quasi-coherent-parasitic-cohomology. Justifies the four-way splitting.
\item Pushforwards of parasitic modules are parasitic. Stacks Project, Cohomology of Algebraic Stacks, Lemma lemma-pushforward-parasitic. Kills the \$R\textasciicircum{}i f\_* \textbackslash{}mathcal\{P\}\_j\$ terms.
\item Kernel and cokernel of \$Q(\textbackslash{}mathcal\{F\}) \textbackslash{}to \textbackslash{}mathcal\{F\}\$ are parasitic. Stacks Project, Cohomology of Algebraic Stacks, Lemma lemma-adjoint-kernel-parasitic. Explains why parasitic terms vanish under \$Q\$.
\item Definition of cohomological \$\textbackslash{}delta\$-functor. Stacks Project, Homology, Definition homology-definition-cohomological-delta-functor. Supplies the target notion.
\end{itemize}\medskip
\ifdefined\hypertarget\hypertarget{sourcepointer-1}{}\fi
\noindent\textit{Source reference.} Original proposition 0771: sources/stacks-chow/tags.txt, source lines 9186--9186. Cross-chapter standard prerequisite cited by the author inside the counted span; it is used as a pointer and its own proof is not part of this item.\par\medskip
\ifdefined\hypertarget\hypertarget{sourcepointer-2}{}\fi
\noindent\textit{Source reference.} Original lemma 0778: sources/stacks-chow/tags.txt, source lines 9193--9193. Cross-chapter standard prerequisite cited by the author inside the counted span; it is used as a pointer and its own proof is not part of this item.\par\medskip
\ifdefined\hypertarget\hypertarget{sourcepointer-3}{}\fi
\noindent\textit{Source reference.} Original reference 010Q: sources/stacks-chow/tags.txt, source lines 1336--1336. Cross-chapter standard prerequisite cited by the author inside the counted span; it is used as a pointer and its own proof is not part of this item.\par\medskip
\ifdefined\hypertarget\hypertarget{sourcepointer-4}{}\fi
\noindent\textit{Source reference.} Original lemma 0776: sources/stacks-chow/tags.txt, source lines 9191--9191. Cross-chapter standard prerequisite cited by the author inside the counted span; it is used as a pointer and its own proof is not part of this item.\par\medskip
\ifdefined\hypertarget\hypertarget{sourcepointer-5}{}\fi
\noindent\textit{Source reference.} Original lemma 0775: sources/stacks-chow/tags.txt, source lines 9190--9190. Cross-chapter standard prerequisite cited by the author inside the counted span; it is used as a pointer and its own proof is not part of this item.\par\medskip
\ifdefined\hypertarget\hypertarget{sourcepointer-6}{}\fi
\noindent\textit{Source reference.} Original lemma 0779: sources/stacks-chow/tags.txt, source lines 9194--9194. Cross-chapter standard prerequisite cited by the author inside the counted span; it is used as a pointer and its own proof is not part of this item.\par\medskip
\begin{proposition}
\label{proposition-direct-image-quasi-coherent}
Let $f : \mathcal{X} \to \mathcal{Y}$ be a quasi-compact and quasi-separated
morphism of algebraic stacks. The functor
$f^* : \QCoh(\mathcal{O}_\mathcal{Y}) \to
\QCoh(\mathcal{O}_\mathcal{X})$
has a right adjoint
$$
f_{\QCoh, *} :
\QCoh(\mathcal{O}_\mathcal{X})
\longrightarrow
\QCoh(\mathcal{O}_\mathcal{Y})
$$
which can be defined as the composition
$$
\QCoh(\mathcal{O}_\mathcal{X}) \to
\textit{LQCoh}^{fbc}(\mathcal{O}_\mathcal{X})
\xrightarrow{f_*} \textit{LQCoh}^{fbc}(\mathcal{O}_\mathcal{Y})
\xrightarrow{Q} \QCoh(\mathcal{O}_\mathcal{Y})
$$
where the functors $f_*$ and $Q$ are as in
Proposition \ref{proposition-loc-qcoh-flat-base-change}
and
Lemma \ref{lemma-adjoint}.
Moreover, if we define $R^if_{\QCoh, *}$ as the composition
$$
\QCoh(\mathcal{O}_\mathcal{X}) \to
\textit{LQCoh}^{fbc}(\mathcal{O}_\mathcal{X})
\xrightarrow{R^if_*} \textit{LQCoh}^{fbc}(\mathcal{O}_\mathcal{Y})
\xrightarrow{Q} \QCoh(\mathcal{O}_\mathcal{Y})
$$
then the sequence of functors $\{R^if_{\QCoh, *}\}_{i \geq 0}$
forms a cohomological $\delta$-functor.
\end{proposition}

\begin{proof}
This is a combination of the results mentioned in the statement.
The adjointness can be shown as follows: Let $\mathcal{F}$
be a quasi-coherent $\mathcal{O}_\mathcal{X}$-module and let
$\mathcal{G}$ be a quasi-coherent $\mathcal{O}_\mathcal{Y}$-module.
Then we have
\begin{align*}
\Mor_{\QCoh(\mathcal{O}_\mathcal{X})}(f^*\mathcal{G}, \mathcal{F})
& =
\Mor_{\textit{LQCoh}^{fbc}(\mathcal{O}_\mathcal{Y})}
(\mathcal{G}, f_*\mathcal{F}) \\
& =
\Mor_{\QCoh(\mathcal{O}_\mathcal{Y})}(\mathcal{G}, Q(f_*\mathcal{F}))
\\
& =
\Mor_{\QCoh(\mathcal{O}_\mathcal{Y})}(\mathcal{G},
f_{\QCoh, *}\mathcal{F})
\end{align*}
the first equality by adjointness of $f_*$ and $f^*$ (for arbitrary sheaves
of modules). By
Proposition \ref{proposition-loc-qcoh-flat-base-change}
we see that $f_*\mathcal{F}$ is an object of
$\textit{LQCoh}^{fbc}(\mathcal{O}_\mathcal{Y})$
(and can be computed in either the fppf or \'etale topology) and we
obtain the second equality by Lemma \ref{lemma-adjoint}. The third
equality is the definition of $f_{\QCoh, *}$.

\medskip\noindent
To see that $\{R^if_{\QCoh, *}\}_{i \geq 0}$ is a cohomological
$\delta$-functor as defined in
Homology, Definition \ref{homology-definition-cohomological-delta-functor}
let
$$
0 \to \mathcal{F}_1 \to \mathcal{F}_2 \to \mathcal{F}_3 \to 0
$$
be a short exact sequence of $\QCoh(\mathcal{O}_\mathcal{X})$.
This sequence may not be an exact sequence in
$\textit{Mod}(\mathcal{O}_\mathcal{X})$ but we know that it is
up to parasitic modules, see
Lemma \ref{lemma-exact-sequence-quasi-coherent-parasitic-cohomology}.
Thus we may break up the sequence into short exact sequences
$$
\begin{matrix}
0 \to \mathcal{P}_1 \to \mathcal{F}_1 \to \mathcal{I}_2 \to 0 \\
0 \to \mathcal{I}_2 \to \mathcal{F}_2 \to \mathcal{Q}_2 \to 0 \\
0 \to \mathcal{P}_2 \to \mathcal{Q}_2 \to \mathcal{I}_3 \to 0 \\
0 \to \mathcal{I}_3 \to \mathcal{F}_3 \to \mathcal{P}_3 \to 0
\end{matrix}
$$
of $\textit{Mod}(\mathcal{O}_\mathcal{X})$ with $\mathcal{P}_i$ parasitic.
Note that each of the sheaves
$\mathcal{P}_j$, $\mathcal{I}_j$, $\mathcal{Q}_j$ is an object of
$\textit{LQCoh}^{fbc}(\mathcal{O}_\mathcal{X})$, see
Proposition \ref{proposition-loc-qcoh-flat-base-change}.
Applying $R^if_*$ we obtain long exact sequences 
$$
\begin{matrix}
0 \to f_*\mathcal{P}_1 \to f_*\mathcal{F}_1 \to f_*\mathcal{I}_2 \to
R^1f_*\mathcal{P}_1 \to \ldots \\
0 \to f_*\mathcal{I}_2 \to f_*\mathcal{F}_2 \to f_*\mathcal{Q}_2 \to
R^1f_*\mathcal{I}_2 \to \ldots \\
0 \to f_*\mathcal{P}_2 \to f_*\mathcal{Q}_2 \to f_*\mathcal{I}_3 \to
R^1f_*\mathcal{P}_2 \to \ldots \\
0 \to f_*\mathcal{I}_3 \to f_*\mathcal{F}_3 \to f_*\mathcal{P}_3 \to
R^1f_*\mathcal{I}_3 \to \ldots
\end{matrix}
$$
where are the terms are objects of
$\textit{LQCoh}^{fbc}(\mathcal{O}_\mathcal{Y})$ by
Proposition \ref{proposition-loc-qcoh-flat-base-change}.
By
Lemma \ref{lemma-pushforward-parasitic}
the sheaves $R^if_*\mathcal{P}_j$ are parasitic, hence vanish on applying
the functor $Q$, see
Lemma \ref{lemma-adjoint-kernel-parasitic}.
Since $Q$ is exact the maps
$$
Q(R^if_*\mathcal{F}_3)
\cong
Q(R^if_*\mathcal{I}_3)
\cong
Q(R^if_*\mathcal{Q}_2)
\rightarrow
Q(R^{i + 1}f_*\mathcal{I}_2)
\cong
Q(R^{i + 1}f_*\mathcal{F}_1)
$$
can serve as the connecting map which turns the family of functors
$\{R^if_{\QCoh, *}\}_{i \geq 0}$
into a cohomological $\delta$-functor.
\end{proof}

\end{document}
